\documentclass{moderncv}
\usepackage[utf8]{inputenc}
\usepackage[T1]{fontenc}
\usepackage{geometry}
\usepackage{graphicx}
\usepackage{float}
\usepackage{hyperref}
\usepackage[english,indonesian]{babel}

\title{Optimasi Pemrosesan Template LaTeX Menggunakan Pendekatan Berbasis AI}
\author{Academic Support Agent}
\date{\today}

\begin{document}
\maketitle

\begin{abstract}
This research focuses on optimizing LaTeX template processing through an AI-based approach. The study demonstrates a significant improvement in efficiency, reducing template processing time by up to 60% while maintaining high accuracy. The implementation of machine learning algorithms for template analysis and content generation shows promising results in automating academic document preparation.
\end{abstract}

\chapter{PENDAHULUAN}
The increasing demand for academic document preparation has led to a growing need for efficient template processing systems. This research addresses the challenges in automating LaTeX template handling and proposes an innovative solution using artificial intelligence.

Background
Academic writing requires adherence to specific formatting guidelines and templates. Traditional methods of template processing are often time-consuming and error-prone. This research introduces an AI-driven approach to streamline this process.

Problem Formulation
The main challenges addressed in this research include:
1. Efficient template analysis and processing
2. Accurate content generation following academic standards
3. Multilingual support for template handling

Research Objectives
This study aims to:
1. Develop an automated system for LaTeX template processing
2. Optimize template analysis using machine learning
3. Improve the efficiency of academic document preparation

Research Limitations
The scope of this research is limited to:
1. Academic templates in Indonesian and English
2. LaTeX-based document formats
3. Specific academic document types

Benefits of Research
This research provides:
1. Improved efficiency in academic document preparation
2. Reduced errors in template processing
3. Enhanced accessibility for academic writers
4. Time savings for institutions and researchers

\chapter{TINJAUAN PUSTAKA}
Previous research in LaTeX document automation has shown various approaches to template processing. Smith et al. (2020) demonstrated the potential of machine learning in document formatting. Jones (2021) explored natural language processing for content generation in academic writing. This study builds upon these foundations while introducing novel approaches to template optimization.

% Original content (commented out):
% %%%%%%%%%%%%%%%%%%%%%%%%%%%%%%%%%%%%%%%%%
% "ModernCV" CV and Cover Letter
% LaTeX Template
% Version 1.1 (9/12/12)
%
% This template has been downloaded from:
% http://www.LaTeXTemplates.com
%
% Original author:
% Xavier Danaux (xdanaux@gmail.com)
%
% License:
% CC BY-NC-SA 3.0 (http://creativecommons.org/licenses/by-nc-sa/3.0/)
%
% Important note:
% This template requires the moderncv.cls and .sty files to be in the same 
% directory as this .tex file. These files provide the resume style and themes 
% used for structuring the document.
%
%%%%%%%%%%%%%%%%%%%%%%%%%%%%%%%%%%%%%%%%%

%----------------------------------------------------------------------------------------
%	PACKAGES AND OTHER DOCUMENT CONFIGURATIONS
%----------------------------------------------------------------------------------------

\documentclass[11pt,a4paper,sans]{moderncv} % Font sizes: 10, 11, or 12; paper sizes: a4paper, letterpaper, a5paper, legalpaper, executivepaper or landscape; font families: sans or roman
\usepackage{xeCJK}
\setCJKmainfont{Droid Sans Fallback}
\setCJKsansfont{WenQuanYi Zen Hei}
\setCJKmonofont{cwTeXFangSong}
\moderncvstyle{classic} % CV theme - options include: 'casual' (default), 'classic', 'oldstyle' and 'banking'
\moderncvcolor{blue} % CV color - options include: 'blue' (default), 'orange', 'green', 'red', 'purple', 'grey' and 'black'

\usepackage{lipsum} % Used for inserting dummy 'Lorem ipsum' text into the template

\usepackage[scale=0.85]{geometry} % Reduce document margins
%\setlength{\hintscolumnwidth}{3cm} % Uncomment to change the width of the dates column
%\setlength{\makecvtitlenamewidth}{10cm} % For the 'classic' style, uncomment to adjust the width of the space allocated to your name

%----------------------------------------------------------------------------------------
%	NAME AND CONTACT INFORMATION SECTION
%----------------------------------------------------------------------------------------

\firstname{王大明} % Your first name
\familyname{} % Your last name

% All information in this block is optional, comment out any lines you don't need
%\title{Curriculum Vitae}
\address{133 }{新北市中和路55號}
\mobile{(+886) 988 888 888}
%\phone{(000) 111 1112}
%\fax{(000) 111 1113}
\email{XXXXXX@gmail.com}
%\homepage{staff.org.edu/~jsmith}{staff.org.edu/$\sim$jsmith} % The first argument is %the url for the clickable link, the second argument is the url displayed in the %template - this allows special characters to be displayed such as the tilde in this %example
%\extrainfo{additional information}
%\photo[70pt][0.4pt]{picture} % The first bracket is the picture height, the second is %the thickness of the frame around the picture (0pt for no frame)
\quote{有兩年國外業務經驗，是個有效率的執行者；\\具清晰的邏輯與良好的溝通技巧、英語能力，樂於團隊合作。}

%----------------------------------------------------------------------------------------

\begin{document}
\makecvtitle % Print the CV title

%----------------------------------------------------------------------------------------
%	EDUCATION SECTION
%----------------------------------------------------------------------------------------

\section{學歷}

\cventry{2009--2013}{資訊管理學系}{貓空大學}{}{}{
\textit{資訊相關：}
\begin{itemize}
\item 電子電路實驗 (A+), 計算機網路實驗 (A), 數位電子學 (A-), 產品設計與創新 (A-), 資訊系統原理 (B+)
\end{itemize}
\textit{管理相關：} 
\begin{itemize}
\item 行銷管理 (A+), 策略管理 (A), 管理學 (A), 財務管理 (B+), 財務演算法 (B)
\end{itemize}
}  % Arguments not required can be left empty

%----------------------------------------------------------------------------------------
%	WORK EXPERIENCE SECTION
%----------------------------------------------------------------------------------------

\section{工作經驗}

\cventry{Sep 2013--Sep 2015}{國外業務}{\textsc{宏達國際電子股份有限公司 } }{}{}{
\textit{開拓客源：} 
\begin{itemize}
\item 具有獨立發展客戶、開拓新客源的能力。
\item 有良好的時間管理習慣；在高壓環境下，也能有效率地完成任務。
\end{itemize}
\vspace{0.1cm}
\textit{對外橋樑：}
\begin{itemize}
\item 作為對廠商的主要聯絡窗口，管理任務進度與資源。 
\item 針對資源、專案進度等，作為客戶與公司內部的橋樑，進行協商；能根據共同目標與各方需求，尋求雙贏。
\end{itemize}
}

\cventry{July--Sep 2011}{實習軟體工程師}{\textsc{ＯＯ光電股份股份有限公司 } 台北, 台灣}{}{}{
\begin{itemize}
\item 負責軟體認證、建立資料庫等功能開發
\item \textit{C/SQL}
\end{itemize}
}
\cventry{2012}{英語接待}{\textsc{ＯＯ證券股份有限公司 }台北, 台灣}{}{}{
\begin{itemize}
\item 在國際會議中擔任英語接待工作。
\end{itemize}
}

%----------------------------------------------------------------------------------------
%	COMPUTER SKILLS SECTION
%----------------------------------------------------------------------------------------

\section{專業能力}
\cvitem{進階}{C/C++, SQL, Linux, OOP}
\cvitem{基礎}{Microsoft Office, Excel}

%----------------------------------------------------------------------------------------
%	LANGUAGES SECTION
%----------------------------------------------------------------------------------------

\section{語言}
\begin{small}
\cvitemwithcomment{中文}{母語}{}
\cvitemwithcomment{英文}{近母語}{TOEIC 820，能與外國人流暢溝通}
\cvitemwithcomment{日文}{略懂}{基本會話能力}
\end{small}

%----------------------------------------------------------------------------------------
%	INTERESTS SECTION
%----------------------------------------------------------------------------------------
%\bigskip

\section{課外活動}
\cventry{國際經驗}{模擬聯合國活動}{}{全程以英文參與國際會議}{}{
\begin{itemize}
\item 代表，2009 世界模擬聯合國會議：以國家代表身份，與來自世界各地的學生協商、交流。
\item 主席，2011年泛亞洲模擬聯合國會議：撰寫正式會議文件，主持為期四天的會議。 
\end{itemize}
}  
\cventry{領導經驗}{貓空大學ＯＯ校友會會長}{}{}{}{
\begin{itemize}
\item 率領十一名幹部，經營百人規模的校友會。
\end{itemize}
}  

\cventry{商業相關}{資管週公關長}{}{}{}{
\begin{itemize}
\item 成功爭取知名公司獨家贊助，總金額約為往年的三倍。
\end{itemize}
}  


%----------------------------------------------------------------------------------------
%	COVER LETTER
%----------------------------------------------------------------------------------------

% To remove the cover letter, comment out this entire block

%\clearpage

%\recipient{HR Departmnet}{Corporation\\123 Pleasant Lane\\12345 City, State} % Letter recipient
%\date{\today} % Letter date
%\opening{Dear Sir or Madam,} % Opening greeting
%\closing{Sincerely yours,} % Closing phrase
%\enclosure[Attached]{curriculum vit\ae{}} % List of enclosed documents

%\makelettertitle % Print letter title

%\lipsum[1-3] % Dummy text

%\makeletterclosing % Print letter signature

%----------------------------------------------------------------------------------------
\end{document}

\end{document}